\section{Capped free-space propagation}
\label{app:repair}
For the dimensional time problem in one dimension, let
$l=\lambda_v$ with $l'=-\lambda_x$ and $\lambda_x$ constant.
On a powered fixed-sign interval, write $|l|=L_a+nt$.
For fixed exhaust $V$ set $F=2P/V$, $\sigma=2P/V^2$ and
$C_m=\lambda_m m^2$. Differentiating the optimized Hamiltonian gives
\begin{equation}
 m'=-\sigma,\quad
 \lambda_m'=\frac{F|l|}{m^2},\quad
 C_m'=\frac{2P|l|}{V}-\frac{4PC_m}{mV^2}.
\end{equation}
Thus $C_m$ is not constant on this boundary. With
$B=n/\sigma$ and $A=L_a+Bm_a$,
\begin{equation}
 \begin{gathered}
 m=m_a-\sigma t,\\
 \lambda_m=\lambda_{m,a}
 +V\left[A\left(\frac1m-\frac1{m_a}\right)
                  +B\log\frac m{m_a}\right].
 \end{gathered}
\end{equation}
For thrust sign $\epsilon=\pm1$,
\begin{equation}
 \begin{split}
 v&=v_a+\epsilon V\log(m_a/m),\\
 x&=x_a+v_at+\epsilon V
                \left[t-\frac m\sigma\log(m_a/m)\right].
 \end{split}
\end{equation}
On an ideal interior arc $C_m$ is constant and $a=Pl/C_m$;
on a coast both $m$ and $\lambda_m$ are constant.

\subsection{Relations carried over from the uncapped analysis}
The pointwise selector holds when evaluated with the current $C_m$.
The fixed-exhaust state maps above, the interior-arc mass formula, and
the exchange argument excluding an avoidable powered--coast rearrangement
under its monotone braking hypotheses also hold. On an interior arc
$l=l_a-nt$,
\begin{equation}
 \frac1{m(t)}=\frac1{m_a}+\frac{P}{2C_m^2}
       (l_a^2t-l_an t^2+n^2t^3/3).
\end{equation}
Integrating $a=Pl/C_m$ gives quadratic velocity and cubic position.
The raw primer remains affine globally; its scaled acceleration does not.
Because $C_m$ varies on boundary arcs, switching times cannot be
obtained from a constant-$C_m$ primer, not even with its terminal value.

The exchange argument uses states only. In a
coast/ceiling-brake mixture with duty fraction $\delta\in[0,1]$, write
$s(t)=\int\delta\,dt$. Then $m=m_a-\sigma s$ and
$v=v_a-V\log(m_a/m)$. Final mass and velocity depend only on total
powered time. Exchanging a burn followed by a coast of length $c$ for
coast followed by burn increases distance by $c\,\Delta v>0$, where
$\Delta v$ is the burn's velocity decrement. Repetition puts all coasting
first. If the braking interval has $0\le v\le v_a$ and $v_a>0$, its
original coast distance is reproduced by a shorter leading coast of
duration $\int_{\rm coast}v\,dt/v_a$. Mass, terminal velocity and distance
are unchanged, while elapsed time strictly decreases whenever a positive
coast follows a positive brake. Approximation extends this state argument
to fractional mixtures.

\subsection{Ceiling-only topology and the unresolved two-sided case}
Consider the ceiling-only, positive-distance, rest-to-rest problem with
$C_m>0$ and positive dry mass. Normalize the constant $\lambda_x=n>0$;
the primer is $l=n(t_z-t)$. Define
$z=Vm|l|/(2C_m)$. The selector gives interior $\mathsf U$ for $z>1$,
ceiling burn $\mathsf B$ for $1/2<z<1$, and coast $\mathsf C$ for $z<1/2$.
On a ceiling arc,
\begin{equation}
 \frac{z'}z=\frac{|l|'}{|l|}+\frac{\sigma}{m}(1-2z),
 \qquad \sigma=2P/V^2.
\end{equation}
On an interior arc $z'/z=|l|'/|l|-\sigma z^2/m$; on a coast it is
$|l|'/|l|$. During boost $|l|'<0$, so $z$ decreases strictly through
$\mathsf U_+,\mathsf B_+,\mathsf C$, with initial truncations allowed.
On braking $|l|'=n>0$. At $z=1/2$ both the coast and boundary vector
fields have $z'>0$. Thus $\mathsf C\to\mathsf B_-$ is regular and cannot
return to the tie: every encounter with that level would point upward.
No sliding arc or terminal $m_6$ event is available.

At $z=1$ both powered fields give
$z'/z=n/|l|-\sigma/m$. The quantity
$k=mn/(\sigma|l|)$ decreases strictly throughout braking, since $m$
decreases and $|l|$ increases. An upward crossing $\mathsf B_-\to\mathsf U_-$
can occur only while $k>1$, and a subsequent downward crossing
$\mathsf U_-\to\mathsf B_-$ only when $k<1$; no further reversal is
possible. Consequently every such finite-switch normal candidate has
the ordered pattern
\begin{equation}
 \mathsf U_+\;\mathsf B_+\;\mathsf C\;\mathsf B_-\;
 \mathsf U_-\;\mathsf B_-,
\end{equation}
with absent intervals and endpoint truncations permitted. In particular
the last boundary burn ends at the endpoint, not at a return to coast.
This is a ceiling-only one-dimensional result. It does not classify the
game's two-sided $50$--$250$ km/s engine. In particular, the $50$ km/s
lower exhaust limit caps acceleration and introduces additional arcs.

To state precisely what remains unresolved here, write $V_-=50$ km/s,
$V_+=250$ km/s, $\rho_V=V_+/V_-=5$, and define $z$ with $V=V_+$.
For $C_m>0$, the dimensional selector becomes
\begin{equation}
 |a|=\begin{cases}
 0,&z<1/2,\\
 2P/(mV_+),&1/2<z<1,\\
 P|l|/C_m,&1<z<\rho_V,\\
 2P/(mV_-),&z>\rho_V.
 \end{cases}
\end{equation}
The last region is a maximum-acceleration ($a_+$) arc, denoted
$\mathsf A_\pm$ according to thrust sign. Its state and mass-adjoint
maps are the fixed-exhaust maps above with $V=V_-$, and
\begin{equation}
 \frac{z'}z=\frac{|l|'}{|l|}
   +\frac{\sigma_-}{m}\left(1-\frac{2z}{\rho_V}\right),
 \qquad\sigma_-=2P/V_-^2.
\end{equation}
Thus the selector and propagation are available for both boundaries,
but a complete two-sided switching-topology theorem, including all
$\mathsf A_\pm$ visits and admissible endpoint truncations, is left open,
as are gravity, abnormal extremals and other endpoint conditions.

\subsection{Endpoint reduction}
This reduction and its numerical example retain the ceiling-only
assumptions of the proved topology.
Use homogeneity to fix $n=1$. The two shooting unknowns are $t_z$ and
$\lambda_m(0)>0$. Propagate the exact interior, boundary and coast maps,
switching with the current $z$, until the first exhaustion $m=m_d$.
Call that time $T(t_z,\lambda_m(0))$. Solve
\begin{equation}
 x(T;t_z,\lambda_m(0))=L,\qquad
 v(T;t_z,\lambda_m(0))=0.
\end{equation}
Reject trajectories without exhaustion or with invalid adjoint sign or
switch inequalities. Alternatively use $T$ as a third unknown and add
$m(T)=m_d$ as a third residual. Normalize the free-time Hamiltonian
afterward if a positive normal scaling exists; and compare competing
solutions, since the reduction does not imply uniqueness.

For $P=1/2$, $V=2$, $m_0=1$, $m_d=1/2$, $L=1$, independent numerical
integration of the selector gives
$t_z=1.22365494$, $\lambda_m(0)=0.61524660$, $T=2.35199587$;
the maximum endpoint residual is $9.6\,10^{-8}$. The coast is approximately
$[0.80144,1.64581]$. Direct transcription with piecewise-constant
acceleration, exact $(x,v,1/m)$ propagation and the ceiling imposed at
each burn cell's final mass gives:
\begin{center}\begin{tabular}{rrr}\toprule
Burn cells&Time&Endpoint error\\\midrule
12&2.35489579&$2.6\,10^{-12}$\\
24&2.35338380&$1.4\,10^{-12}$\\
48&2.35267314&$1.3\,10^{-12}$\\\bottomrule
\end{tabular}\end{center}
The within-cell ceiling holds because inverse mass increases monotonically
and the acceleration is constant. Ceiling violations at optimizer tolerance
are below $9\,10^{-13}$. The transcription independently checks the
trajectory equations and approaches the shooting time from above.

\section{The opposing CW strategy}
\label{app:game}
For $b_P>b_E$ let $H$ solve the game equation and put
$k=\sqrt{b_E}/(\sqrt{b_P}-\sqrt{b_E})$.
The evader starts with $u_E=ka_H$.
Any rendezvous at $t<H$ needs relative correction $a_t$.
Minimum-norm orthogonality and
$\int_0^t a_H\cdot a_t\,ds=C(H)$ imply
\begin{equation}
 \int_0^t(|u_P|^2-|u_E|^2)ds
 \ge C(t)+2kC(H)>b_P(0)-b_E(0).
\end{equation}
Rendezvous cannot occur before, or at, the first resource parity.
At parity the evader switches to the following feedback, which also
handles an initial $b_E\ge b_P$.

Reset the clock at a nonzero state $z_*$ and let
$l=z_*/|z_*|$. Introduce virtual evader reserve
$q(0)=b_P(0)$, $q'=-|u_E|^2$ and define
\begin{equation}
 \begin{gathered}
 Y=l^T\Phi(-t)z,\quad
 c=B^T\Phi(-t)^Tl,\quad \delta=q-b_P,\\
 u_E=(\delta/Y)c.
 \end{gathered}
\end{equation}
Initially $Y>0$, $\delta=0$. While $Y>0$,
$\delta'=|u_P|^2-|u_E|^2$, so $\delta\ge0$, and
\begin{equation}
 (\delta Y^2)'=Y^2|u_P-u_E|^2\ge0.
\end{equation}
If the product becomes positive, $\delta\le q(0)$ prevents any
finite zero of $Y$. If it stays zero, both controls vanish almost
everywhere and $Y$ stays constant. Thus rendezvous is prevented
and the virtual budget bounds actual evader expenditure.
The zero-reserve case is coasting. This completes the opposing
bound under the theorem's information convention.

\section{An explicit nonlinear capture tube}
\label{app:tube}
To make \eqref{eq:cwcert} constructive, define
\begin{equation}
 \begin{split}
 g_h&=h+h^2+\tfrac23h^3,\\
 G_1&=\tfrac12h^2+\tfrac13h^3+\tfrac16h^4,\\
 G_2&=\tfrac13h^3+\tfrac12h^4+\tfrac7{15}h^5
                     +\tfrac29h^6+\tfrac4{63}h^7.
 \end{split}
\end{equation}
For rotating state $s=(q,v)$ put
\begin{equation}
 \begin{split}
 F_h(s)&=|q|+(\tfrac32h^2+h^3)|q_x|+g_h|v|,\\
 D_h&=F_h(z_0)+\sqrt{C(h)G_2},\\
 E_h&=F_h(s_{E0})+\sqrt{b_EG_2}.
 \end{split}
\end{equation}
The two algebraic conditions
\begin{equation}
 \begin{gathered}
 E_h+D_h+N_\epsilon G_1<\epsilon,\\
 \sqrt{b_P}\ge\sqrt{b_E}+\sqrt{C(h)}
                      +L_\epsilon\sqrt h\,D_h
 \end{gathered}
\end{equation}
certify nonlinear rendezvous by $h$.
Indeed $\norm{\Pi\Phi(t)B}\le t+t^2+2t^3/3$ and
$|\Pi\Phi(t)s|\le F_h(s)$ for $0\le t\le h$.
Variation of constants bounds the evader by
$E_h+N_\epsilon G_1$ and the relative path by $D_h$.
The first inequality prevents a first exit of either ship.
The second bounds the compensation effort.
For a physical exclusion radius also require
$R_0(1-\epsilon)\ge r_{\min}$.

Failure of either inequality is inconclusive.
Appendix~\ref{app:inertial} supplies an inertial half-space certificate for wider separations.

\section{An inertial half-space rendezvous certificate}
\label{app:inertial}
Here all quantities are physical and $g(x)=-\mu x/|x|^3$; use $\mu=1$
for normalized quantities. Put $r_0=x_{E0}-x_{P0}$,
$w_0=v_{E0}-v_{P0}$ and $s=t/h$. The free-space rendezvous correction
and its relative path are
\begin{align}
 a_h(t)&=6r_0/h^2+4w_0/h-(12r_0/h^3+6w_0/h^2)t,\\
 r_J(t)&=(1-3s^2+2s^3)r_0+h(s-2s^2+s^3)w_0.
\end{align}
Thus $\int|a_h|^2=C_J(h)$ from Section~\ref{sec:interception}, and direct
polynomial integration gives
\begin{equation}
 I_J=\int_0^h|r_J|^2dt=
 \frac{13h}{35}|r_0|^2+\frac{11h^2}{105}r_0\cdot w_0
                      +\frac{h^3}{105}|w_0|^2.
\end{equation}
Choose a unit vector $n$ and $r_*>0$. Require, for every $0\le t\le h$,
\begin{equation}
\begin{split}
 n\cdot(x_{E0}+tv_{E0})-\frac{\mu t^2}{2r_*^2}
 -\sqrt{b_Et^3/3}\\
 -\max\{n\cdot r_J(t),0\}>r_*.
\end{split}
\label{eq:halfspace}
\end{equation}
A simpler sufficient scalar check replaces the linear term by the minimum
of its two endpoint values and subtracts
$\mu h^2/(2r_*^2)+\sqrt{b_Eh^3/3}+D_J$, where
$D_J=|r_0|+4h|w_0|/27$.
If also
\begin{equation}
 \sqrt{b_P}\ge\sqrt{b_E}+\sqrt{C_J(h)}
                   +\frac{2\mu}{r_*^3}\sqrt{I_J},
\end{equation}
the response $a_P=a_E+g(x_E)-g(x_P)+a_h$ guarantees full-state
rendezvous by $h$. To prove the corridor, before a first exit the evader
gravitational displacement in direction $n$ is bounded below by
$-\mu t^2/(2r_*^2)$ and its controlled displacement by
$-\sqrt{b_Et^3/3}$. Since $x_P=x_E-r_J$, \eqref{eq:halfspace} excludes
a first exit of either ship. Their joining segment is in the convex
half-space $n\cdot x\ge r_*$, so
$|g(x_E)-g(x_P)|\le K_g|r_J|$, $K_g=2\mu/r_*^3$.
Minkowski and the quadrature then prove the resource bound. Choose
$r_*\ge r_{\min}$ if an exclusion floor is imposed. A failed sufficient
test is inconclusive, and finite-band admissibility is an additional check.

\section{CW entries and the rational Jensen calculation}
\label{app:interval}
Write $S=\sin h$, $C_h=\cos h$, $S_2=\sin(2h)$. The unit-mean-motion
transition matrix is
\begin{equation}
\Phi(h)=\begin{pmatrix}
4-3C_h&0&S&2(1-C_h)\\
6(S-h)&1&-2(1-C_h)&4S-3h\\
3S&0&C_h&2S\\
-6(1-C_h)&0&-2S&4C_h-3
\end{pmatrix}.
\end{equation}
The symmetric Gramian has upper-triangle entries
\begin{align*}
W_{11}&=13h/2-8S+3S_2/4,\\
W_{12}&=-3(h-S)^2,\\
W_{13}&=4(1-C_h)-3S^2/2,\\
W_{14}&=-11h+14S-3S_2/2,\\
W_{22}&=3h^3+14h+24hC_h-32S-3S_2,\\
W_{23}&=5h+6hC_h-8S-3S_2/2,\\
W_{24}&=6S^2+4(1-C_h)-12hS+9h^2/2,\\
W_{33}&=5h/2-3S_2/4,\\
W_{34}&=3S^2+6(C_h-1),\\
W_{44}&=19h+3S_2-24S.
\end{align*}
These follow by integrating $\Phi(t)BB^T\Phi(t)^T$ entrywise.
For the Jensen example, $y=\Phi(h)\xi=(1,1-3h/2,0,-3/2)^T$.
Evaluate $c=y^TW^{-1}y$ at rational $h=10,21/2,11$ using interval
Gaussian elimination. The following construction specifies the interval
calculation without requiring an external supplement. For each sine or
cosine at argument $x\in\{h,2h\}$, sum its rational Taylor polynomial
through degree 181 or 180, respectively; append the symmetric Lagrange
remainder interval $[-|x|^{d+1}/(d+1)!,|x|^{d+1}/(d+1)!]$.
Convert fractions outward to 70 decimal digits. Round every subsequent
interval arithmetic operation outward, check that each elimination pivot
excludes zero, solve $Wq=y$, and form $y^Tq$.
The deliberately widened enclosures are $|c(h)-\bar c(h)|\le2\,10^{-20}$,
with the following rational centers:
\begin{center}\begin{tabular}{rl}\toprule
$h$&$\bar c(h)$\\\midrule
10&0.08299553059490180337\\
10.5&0.08164450969532354695\\
11&0.08014088306972455432\\\bottomrule
\end{tabular}\end{center}
Consequently the Jensen gap differs from
$7.630286301036810\,10^{-5}$ by at most $4\,10^{-20}$ and is strictly positive.
The argument in Section~\ref{sec:timing} transfers this leading sign to
the nonlinear value for sufficiently small perturbations; it does not
supply an explicit largest perturbation.
